\section{The routes}\label{sec:routes}

We now treat the statements as propositional letters and the results of
the paper as rules between them. Besides the seven statements of
\Cref{def:statements} and \HC{} there are three more: \st{SR}, that the
subvariety $Y$ of \Cref{thm:E} is semiregular for some $(X,L)$ of each genus
$n+1\ge5$ with maximal Hodge group; \st{W3}, that the Weil classes are
algebraic on every abelian variety of split Weil type for $\QQ(\sqrt{-3})$;
and \st{Fer}, the Hodge conjecture for every Fermat variety. The rules are
listed in \Cref{tab:rules}.

\begin{table}[ht]
\centering
\renewcommand{\arraystretch}{1.18}
\begin{tabular}{@{}l>{\raggedright\arraybackslash}p{0.56\textwidth}@{}}
\toprule
Implication & Source\\
\midrule
$\st{F2}\wedge\st{F3$'$}\Rightarrow\HC$ & \Cref{ssec:proofA}: if
\st{F2} holds then $\Ab^{p}(X)=\Alg^{p}(X)$\\
$\HC\Rightarrow\st{F2},\st{F3$'$},\st{L},\st{M},\st{V}$ &
\Cref{lem:hcimplies}\\
$\st{L}\wedge\st{M}\Rightarrow\HC$ & \Cref{prop:LM}, after Andr\'e
\cite{And96}\\
$\st{L}\Rightarrow\st{F2}$ & Abdulali, Andr\'e \cite[Theorem~4]{Mil20}\\
$\st{V}\Rightarrow\st{F2}$ & \cite[Theorem~3 and Remark~2]{Mil20}\\
$\st{V}\Rightarrow\st{IP}$ & \st{IP} is a case of \st{V}\\
$\st{F3$'$}\Rightarrow\st{M}$ & \Cref{prop:f3m}\\
$\st{CM}\wedge\st{IP}\Rightarrow\st{F2}$ & \Cref{prop:cmip}\\
$\st{F2}\Rightarrow\st{CM},\st{IP}$ & \Cref{prop:cmip}\\
$\st{CM}\Rightarrow\st{Fer}$, $\HC\Rightarrow\st{Fer}$ & \Cref{thm:C}\\
$\st{SR}\Rightarrow\st{W3}$ & \Cref{cor:schoensemireg} for $2n\ge8$;
\Cref{thm:lef11} and \cite{Sch88,Sch98} for $2n\le6$\\
$\st{F2}\Rightarrow\st{W3}$ & \st{W3} is a case of \st{F2}\\
\bottomrule
\end{tabular}
\caption{The implications between the statements. A set of statements
\emph{implies} a statement if the statement is reached from the set by
repeated use of these rules.}
\label{tab:rules}
\end{table}

For $2n\le6$ the conclusion of \st{SR} is already known, which is why
\st{SR} only concerns the genera $n+1\ge5$. A set of
statements is \emph{closed} if it contains the conclusion of every rule whose
premises it contains. The statements implied by a set $S$ form the smallest
closed set containing $S$, so a closed set that contains $S$ and not \HC{}
shows that \HC{} does not follow from $S$.

\begin{lemma}\label{lem:closedsets}
Let $\mathcal{U}$ be the set of all eleven statements. The four sets
\[
\begin{aligned}
  T_{1}&=\mathcal{U}\setminus\{\HC,\st{M},\st{F3$'$}\},&
  T_{2}&=\mathcal{U}\setminus\{\HC,\st{L},\st{F3$'$}\},\\
  T_{3}&=\mathcal{U}\setminus\{\HC,\st{F2},\st{L},\st{V},\st{IP}\},&
  T_{4}&=\mathcal{U}\setminus\{\HC,\st{F2},\st{L},\st{V},\st{CM}\}
\end{aligned}
\]
are closed, and so is $\{\st{CM},\st{Fer}\}$.
\end{lemma}

\begin{proof}
A rule can fail to be respected by a set only if its conclusion is missing
from the set while all its premises are present. We check the rules whose
conclusions are missing. In $T_{1}$ the missing conclusions are \HC, \st{M}
and \st{F3$'$}; every rule concluding one of them has a premise among \HC,
\st{F3$'$} and \st{M}, which are missing. In $T_{2}$ the rules concluding
\HC, \st{L} or \st{F3$'$} need one of \st{F3$'$}, \HC{} or \st{L}. In
$T_{3}$ every rule concluding \HC, \st{F2}, \st{L}, \st{V} or \st{IP} has a
premise among \st{F2}, \HC, \st{L}, \st{V} and \st{IP}; the rule
$\st{CM}\wedge\st{IP}\Rightarrow\st{F2}$ is blocked by \st{IP}. $T_{4}$ is
the same with \st{CM} in place of \st{IP}. In $\{\st{CM},\st{Fer}\}$ the
only rule with all premises present is $\st{CM}\Rightarrow\st{Fer}$.
\end{proof}

\begin{proof}[Proof of \Cref{thm:B}]
Let $S$ be a set of statements, not containing \HC, that implies \HC. Then
$S$ is contained in none of $T_{1},\dots,T_{4}$, so it meets each of
\[
  \{\st{M},\st{F3$'$}\},\quad\{\st{L},\st{F3$'$}\},\quad
  \{\st{F2},\st{L},\st{V},\st{IP}\},\quad\{\st{F2},\st{L},\st{V},\st{CM}\}.
\]
If $\st{F3$'$}\notin S$, the first two give $\st{M},\st{L}\in S$, which is
(b). If $\st{F3$'$}\in S$ and $S$ contains none of \st{F2}, \st{L}, \st{V},
the last two give $\st{IP},\st{CM}\in S$. This is (a).

Conversely, $\{\st{F2},\st{F3$'$}\}$ and $\{\st{L},\st{M}\}$ imply \HC{}
directly, and $\{\st{L},\st{F3$'$}\}$, $\{\st{V},\st{F3$'$}\}$ and
$\{\st{CM},\st{IP},\st{F3$'$}\}$ imply \st{F2} and then \HC. From \HC{} the
rules give \st{F2}, \st{F3$'$}, \st{L}, \st{M}, \st{V}, then \st{CM},
\st{IP}, \st{W3} and \st{Fer}: every statement except \st{SR}. No minimal set
contains \st{SR}, because \st{SR} is a premise only of the rule that
concludes \st{W3}, and \st{W3} is a premise of no rule.
\end{proof}

\begin{corollary}\label{cor:bypass}
\begin{enumerate}[label=\textup{(\roman*)}]
\item A set of statements that implies \HC{} and contains neither \st{F2}
nor \st{F3$'$} contains \st{L} and \st{M}; by \Cref{thm:A} this pair is
equivalent to \HC.
\item Every set that implies \HC{} also implies \st{F2} and \st{F3$'$}.
\item A set that implies \HC{} still implies it when \st{SR} is removed.
\item \st{CM} implies \st{Fer} and nothing else; neither
$\{\st{CM},\st{F3$'$}\}$ nor $\{\st{CM},\st{SR}\}$ implies \HC.
\end{enumerate}
\end{corollary}

\begin{proof}
(i) is \Cref{thm:B}(b), since (a) requires \st{F3$'$}. (ii) holds because
\HC{} implies \st{F2} and \st{F3$'$}. (iii) is the last sentence of the proof
of \Cref{thm:B}. For (iv), $\{\st{CM},\st{Fer}\}$ is closed by
\Cref{lem:closedsets}, and $\{\st{CM},\st{F3$'$}\}$ and $\{\st{CM},\st{SR}\}$
lie in the closed set $T_{3}$, which does not contain \HC.
\end{proof}

So the conjecture can be reached without assuming \st{F2} or \st{F3$'$},
but only through \st{L} and \st{M}, which carry the whole conjecture between
them; and it cannot be reached without proving \st{F2} and \st{F3$'$} along
the way. Schoen's subvariety, if it is semiregular, settles a family of Weil
classes in every dimension, but it lies on no route to \HC. The closed set
$\{\st{CM},\st{Fer}\}$ is a valuation of the eleven statements in which
every rule of \Cref{tab:rules} holds, \st{CM} holds, and \HC{} fails: the
rules alone cannot derive the conjecture from the preprint \cite{OAI26}.
